\PassOptionsToPackage{unicode}{hyperref}
\PassOptionsToPackage{hyphens}{url}
\PassOptionsToPackage{dvipsnames,svgnames,x11names}{xcolor}
%
\documentclass[
  11pt,
  letterpaper,
  DIV=11,
  numbers=noendperiod]{scrartcl}
\usepackage{xcolor}
\usepackage[margin=1in]{geometry}
\usepackage{amsmath,amssymb}
\setcounter{secnumdepth}{-\maxdimen} % remove section numbering
\usepackage{iftex}
\ifPDFTeX
  \usepackage[T1]{fontenc}
  \usepackage[utf8]{inputenc}
  \usepackage{textcomp}
\else
  \usepackage{unicode-math}
  \defaultfontfeatures{Scale=MatchLowercase}
  \defaultfontfeatures[\rmfamily]{Ligatures=TeX,Scale=1}
\fi
\usepackage{lmodern}
\IfFileExists{upquote.sty}{\usepackage{upquote}}{}
\IfFileExists{microtype.sty}{%
  \usepackage[]{microtype}
  \UseMicrotypeSet[protrusion]{basicmath}
}{}
\makeatletter
\@ifundefined{KOMAClassName}{%
  \IfFileExists{parskip.sty}{%
    \usepackage{parskip}
  }{%
    \setlength{\parindent}{0pt}
    \setlength{\parskip}{6pt plus 2pt minus 1pt}}
}{%
  \KOMAoptions{parskip=half}}
\makeatother

% Code listings (verbatim with light shading)
\usepackage{color}
\usepackage{fancyvrb}
\newcommand{\VerbBar}{|}
\newcommand{\VERB}{\Verb[commandchars=\\\{\}]}
\DefineVerbatimEnvironment{Highlighting}{Verbatim}{commandchars=\\\{\}}
\usepackage{framed}
\definecolor{shadecolor}{RGB}{241,243,245}
\newenvironment{Shaded}{\begin{snugshade}}{\end{snugshade}}
\newcommand{\AlertTok}[1]{\textcolor[rgb]{0.68,0.00,0.00}{#1}}
\newcommand{\AnnotationTok}[1]{\textcolor[rgb]{0.37,0.37,0.37}{#1}}
\newcommand{\AttributeTok}[1]{\textcolor[rgb]{0.40,0.45,0.13}{#1}}
\newcommand{\BaseNTok}[1]{\textcolor[rgb]{0.68,0.00,0.00}{#1}}
\newcommand{\BuiltInTok}[1]{\textcolor[rgb]{0.00,0.23,0.31}{#1}}
\newcommand{\CharTok}[1]{\textcolor[rgb]{0.13,0.47,0.30}{#1}}
\newcommand{\CommentTok}[1]{\textcolor[rgb]{0.37,0.37,0.37}{#1}}
\newcommand{\CommentVarTok}[1]{\textcolor[rgb]{0.37,0.37,0.37}{\textit{#1}}}
\newcommand{\ConstantTok}[1]{\textcolor[rgb]{0.56,0.35,0.01}{#1}}
\newcommand{\ControlFlowTok}[1]{\textcolor[rgb]{0.00,0.23,0.31}{\textbf{#1}}}
\newcommand{\DataTypeTok}[1]{\textcolor[rgb]{0.68,0.00,0.00}{#1}}
\newcommand{\DecValTok}[1]{\textcolor[rgb]{0.68,0.00,0.00}{#1}}
\newcommand{\DocumentationTok}[1]{\textcolor[rgb]{0.37,0.37,0.37}{\textit{#1}}}
\newcommand{\ErrorTok}[1]{\textcolor[rgb]{0.68,0.00,0.00}{#1}}
\newcommand{\ExtensionTok}[1]{\textcolor[rgb]{0.00,0.23,0.31}{#1}}
\newcommand{\FloatTok}[1]{\textcolor[rgb]{0.68,0.00,0.00}{#1}}
\newcommand{\FunctionTok}[1]{\textcolor[rgb]{0.28,0.35,0.67}{#1}}
\newcommand{\ImportTok}[1]{\textcolor[rgb]{0.00,0.46,0.62}{#1}}
\newcommand{\InformationTok}[1]{\textcolor[rgb]{0.37,0.37,0.37}{#1}}
\newcommand{\KeywordTok}[1]{\textcolor[rgb]{0.00,0.23,0.31}{\textbf{#1}}}
\newcommand{\NormalTok}[1]{\textcolor[rgb]{0.00,0.23,0.31}{#1}}
\newcommand{\OperatorTok}[1]{\textcolor[rgb]{0.37,0.37,0.37}{#1}}
\newcommand{\OtherTok}[1]{\textcolor[rgb]{0.00,0.23,0.31}{#1}}
\newcommand{\PreprocessorTok}[1]{\textcolor[rgb]{0.68,0.00,0.00}{#1}}
\newcommand{\RegionMarkerTok}[1]{\textcolor[rgb]{0.00,0.23,0.31}{#1}}
\newcommand{\SpecialCharTok}[1]{\textcolor[rgb]{0.37,0.37,0.37}{#1}}
\newcommand{\SpecialStringTok}[1]{\textcolor[rgb]{0.13,0.47,0.30}{#1}}
\newcommand{\StringTok}[1]{\textcolor[rgb]{0.13,0.47,0.30}{#1}}
\newcommand{\VariableTok}[1]{\textcolor[rgb]{0.07,0.07,0.07}{#1}}
\newcommand{\VerbatimStringTok}[1]{\textcolor[rgb]{0.13,0.47,0.30}{#1}}
\newcommand{\WarningTok}[1]{\textcolor[rgb]{0.37,0.37,0.37}{\textit{#1}}}

% Tables (for the self-assessment guide)
\usepackage{longtable,booktabs,array}

% Callout boxes (note + tip)
\usepackage[skins,breakable]{tcolorbox}
\usepackage{fontawesome5}
\definecolor{quarto-callout-note-color}{HTML}{0758E5}
\definecolor{quarto-callout-note-color-frame}{HTML}{4582ec}
\definecolor{quarto-callout-tip-color}{HTML}{00A047}
\definecolor{quarto-callout-tip-color-frame}{HTML}{02b875}

\setlength{\emergencystretch}{3em}
\providecommand{\tightlist}{%
  \setlength{\itemsep}{0pt}\setlength{\parskip}{0pt}}

\usepackage{bookmark}
\IfFileExists{xurl.sty}{\usepackage{xurl}}{}
\urlstyle{same}
\hypersetup{
  pdftitle={Quiz 1 --- Practice Answer Key},
  colorlinks=true,
  linkcolor={blue},
  filecolor={Maroon},
  citecolor={Blue},
  urlcolor={Blue},
  pdfcreator={LaTeX}}

\title{Quiz 1 --- Practice Answer Key}
\usepackage{etoolbox}
\makeatletter
\providecommand{\subtitle}[1]{%
  \apptocmd{\@title}{\par {\large #1 \par}}{}{}
}
\makeatother
\subtitle{STAT 184 --- Summer 2026}
\author{}
\date{}

\begin{document}
\maketitle

\textbf{Total: 30 points} $\cdot$ 9 questions $\cdot$ 20 minutes

\begin{tcolorbox}[enhanced jigsaw, arc=.35mm, bottomrule=.15mm, bottomtitle=1mm, breakable, colback=white, colbacktitle=quarto-callout-tip-color!10!white, colframe=quarto-callout-tip-color-frame, coltitle=black, left=2mm, leftrule=.75mm, opacityback=0, opacitybacktitle=0.6, rightrule=.15mm, title=\textcolor{quarto-callout-tip-color}{\faLightbulb[regular]}\hspace{0.5em}{How to use this key}, titlerule=0mm, toprule=.15mm, toptitle=1mm]

\begin{itemize}
\tightlist
\item Try the practice quiz first \textbf{without} looking at the key, ideally in a 20-minute closed-book session.
\item Then check your answers below.
\item If you missed a question, the ``Why'' notes will point you to what concept to review.
\end{itemize}

\end{tcolorbox}

\vspace{0.5em}

\hrulefill

\subsection{1. (2 pts) Assignment}

\textbf{Answer: C} (\texttt{score\ \textless{}-\ 100})

\textbf{Why:} \texttt{\textless{}-} is the assignment operator we use in this class. \texttt{=} works in some places but causes confusion later. \texttt{==} is for \textit{checking equality}, not assigning. \texttt{100\ \textless{}-\ score} runs in the wrong direction. \texttt{assign()} exists but isn't taught.

\hrulefill

\subsection{2. (2 pts) Packages}

\textbf{Answer: B} (\texttt{library(ggplot2)})

\textbf{Why:} Installing is a \textbf{one-time} action per computer. Loading with \texttt{library()} happens \textbf{every session} --- packages don't auto-load when R starts.

\hrulefill

\subsection{3. (4 pts, 1 pt each) Predict the output}

\begin{enumerate}
\def\labelenumi{(\alph{enumi})}
\item \texttt{sqrt(144)} $\rightarrow$ \textbf{\texttt{12}}
\item \texttt{c(10,\ 20,\ 30)\ +\ 5} $\rightarrow$ \textbf{\texttt{15\ 25\ 35}} (or \texttt{[1]\ 15\ 25\ 35})
\end{enumerate}

\begin{quote}
Why: scalar arithmetic on a vector is applied \textbf{element-by-element}.
\end{quote}

\begin{enumerate}
\def\labelenumi{(\alph{enumi})}
\setcounter{enumi}{2}
\item \texttt{15\ \%\%\ 4} $\rightarrow$ \textbf{\texttt{3}}
\end{enumerate}

\begin{quote}
Why: \texttt{\%\%} is the \textbf{modulo} operator --- the remainder after division. 15 $\div$ 4 = 3 remainder 3.
\end{quote}

\begin{enumerate}
\def\labelenumi{(\alph{enumi})}
\setcounter{enumi}{3}
\item \texttt{mean(c(1,\ 3,\ 5,\ 7,\ 9))} $\rightarrow$ \textbf{\texttt{5}}
\end{enumerate}

\begin{quote}
Why: (1+3+5+7+9) / 5 = 25 / 5 = 5.
\end{quote}

\hrulefill

\subsection{4. (4 pts, 1 pt each) Identify the parts}

\begin{enumerate}
\def\labelenumi{(\alph{enumi})}
\item Data table: \textbf{\texttt{Movies}}
\end{enumerate}

\begin{quote}
Why: starts with a capital letter, no parentheses after it.
\end{quote}

\begin{enumerate}
\def\labelenumi{(\alph{enumi})}
\setcounter{enumi}{1}
\item Functions: \textbf{any two of} \texttt{filter}, \texttt{summarise}, \texttt{mean}
\end{enumerate}

\begin{quote}
Why: each is followed by an open parenthesis.
\end{quote}

\begin{enumerate}
\def\labelenumi{(\alph{enumi})}
\setcounter{enumi}{2}
\item Variable: \textbf{\texttt{year}} or \textbf{\texttt{rating}}
\end{enumerate}

\begin{quote}
Why: lowercase, used inside function arguments.
\end{quote}

\begin{enumerate}
\def\labelenumi{(\alph{enumi})}
\setcounter{enumi}{3}
\item Constant: \textbf{\texttt{2010}}
\end{enumerate}

\begin{quote}
Why: a literal numeric value.
\end{quote}

\hrulefill

\subsection{5. (3 pts) Data structures}

\textbf{Answer: A}

A one-dimensional collection of values, all of the same type.

\textbf{Why each other option is wrong:}

\begin{itemize}
\tightlist
\item B describes a \textbf{matrix}.
\item C describes a \textbf{data frame}.
\item D describes a \textbf{list}.
\item E describes a single value, not a vector. (Technically a length-1 vector exists, but A is the better description.)
\end{itemize}

\hrulefill

\subsection{6. (3 pts, 1 pt each) Function arguments}

\begin{enumerate}
\def\labelenumi{(\alph{enumi})}
\item Function name: \textbf{\texttt{round}}
\item Number of arguments: \textbf{2} (\texttt{3.14159} and \texttt{digits\ =\ 3})
\item The first argument: \textbf{positional} (it has no name on the left side of \texttt{=})
\end{enumerate}

\begin{quote}
Why: in \texttt{round(3.14159,\ digits\ =\ 3)}, the first argument is matched by \textit{position} --- \texttt{round()} expects \texttt{x} first. The second is matched by \textit{name} (\texttt{digits\ =\ 3}).
\end{quote}

\hrulefill

\subsection{7. (3 pts) Storing results}

\textbf{Full credit (3 pts):}

Both lines compute the mean (which is \texttt{4}), but:

\begin{itemize}
\tightlist
\item \texttt{mean(c(2,\ 4,\ 6))} just \textbf{prints} the result to the console --- it's not saved anywhere.
\item \texttt{m\ \textless{}-\ mean(c(2,\ 4,\ 6))} \textbf{assigns} the result to an object called \texttt{m}, so you can use \texttt{m} later.
\end{itemize}

\textbf{Why:} without \texttt{\textless{}-}, the value is computed and then discarded. The assignment operator is what stores the result for future use.

\hrulefill

\subsection{8. (4 pts, 1 pt each) Histograms}

\begin{enumerate}
\def\labelenumi{(\alph{enumi})}
\item X-axis variable: \textbf{\texttt{bill\_length\_mm}}
\item Bin width: \textbf{2} (each bar covers 2 mm of bill length)
\item Inside color: \textbf{steelblue}
\item Outline color: \textbf{white}
\end{enumerate}

\begin{quote}
Why \texttt{fill} vs.~\texttt{color}: in \texttt{ggplot2}, \texttt{fill} is the inside of a shape and \texttt{color} is the border/outline.
\end{quote}

\hrulefill

\subsection{9. (5 pts) Debug the code}

\begin{enumerate}
\def\labelenumi{(\alph{enumi})}
\tightlist
\item \textbf{(2 pts)} The code uses \texttt{\%\textgreater{}\%} (or \texttt{|\textgreater{}}) between the \texttt{ggplot()} call and \texttt{geom\_histogram()}, but \texttt{ggplot2} layers must be joined with \texttt{+}, not the pipe.
\end{enumerate}

\begin{quote}
Why: the pipe \texttt{|\textgreater{}} (or \texttt{\%\textgreater{}\%}) is for \textbf{data verbs} like dplyr functions. \texttt{ggplot2} was built before the pipe became popular and uses \texttt{+} to add plot layers.
\end{quote}

\begin{enumerate}
\def\labelenumi{(\alph{enumi})}
\setcounter{enumi}{1}
\tightlist
\item \textbf{(3 pts)} Correct version:
\end{enumerate}

\begin{Shaded}
\begin{Highlighting}[]
\FunctionTok{ggplot}\NormalTok{(penguins, }\FunctionTok{aes}\NormalTok{(}\AttributeTok{x =}\NormalTok{ body\_mass\_g)) }\SpecialCharTok{+}
  \FunctionTok{geom\_histogram}\NormalTok{()}
\end{Highlighting}
\end{Shaded}

\begin{quote}
Grading notes:

\begin{itemize}
\tightlist
\item Full credit if the \texttt{+} replaces the \texttt{\%\textgreater{}\%} and the rest is unchanged.
\item 2 pts if they wrote correct code but missed an explanation in (a).
\item 1 pt if they tried but produced something that still wouldn't run.
\end{itemize}
\end{quote}

\hrulefill

\subsection{How did you do?}

\begin{longtable}[]{@{}lp{4.5in}@{}}
\toprule\noalign{}
\textbf{Score} & \textbf{Interpretation} \\
\midrule\noalign{}
\endhead
27--30 & You're ready for the real quiz. \\
22--26 & Solid --- review the questions you missed. \\
17--21 & Worth a serious review. Re-read DC §3 and revisit your Day 2 and Day 3 notes. \\
Below 17 & Come to office hours before the real quiz. We can work through the trouble spots together. \\
\bottomrule\noalign{}
\end{longtable}

\vspace{0.5em}

\begin{tcolorbox}[enhanced jigsaw, arc=.35mm, bottomrule=.15mm, bottomtitle=1mm, breakable, colback=white, colbacktitle=quarto-callout-note-color!10!white, colframe=quarto-callout-note-color-frame, coltitle=black, left=2mm, leftrule=.75mm, opacityback=0, opacitybacktitle=0.6, rightrule=.15mm, title=\textcolor{quarto-callout-note-color}{\faInfo}\hspace{0.5em}{Reminder}, titlerule=0mm, toprule=.15mm, toptitle=1mm]

The real Quiz 1 will be \textbf{closed-book, closed-electronics, no AI}, and worth \textbf{3\% of your final grade}. The questions will be different but cover the same topics.

\end{tcolorbox}

\end{document}
